\documentclass[11pt,a4paper]{article}
\usepackage[margin=25mm]{geometry}
\usepackage[hidelinks]{hyperref}
\usepackage{parskip}
\pagestyle{empty}

\begin{document}

\noindent
Tze Hui Liew (corresponding author)\\
Faculty of Information Science and Technology, Multimedia University\\
Melaka, Malaysia\\
\href{mailto:thliew@mmu.edu.my}{thliew@mmu.edu.my}

\medskip
\noindent\today

\medskip
\noindent
The Editor-in-Chief\\
\emph{Computers \& Security}

\medskip
\noindent
\textbf{Submission of a Full Length Article:} ``It Is Not the Signature: A Pre-Authentication Key Parse Dominates JWT
Validation in jsonwebtoken and Is Load-Bearing for Its Key-Confusion Defence''

Dear Editor,

We submit the above manuscript for consideration as a Full Length Article in
\emph{Computers \& Security}.

JSON Web Token validation is the authentication step of most microservice
requests, and its cost is usually attributed to the signature. We show that in
\texttt{jsonwebtoken}, a widely used Node.js implementation, the signature is
under a tenth of the validation cost. The dominant component is a discarded
attempt to parse the shared secret as an asymmetric key. The manuscript's
central contribution for this journal is the security analysis of that
component:
\begin{itemize}
\item It runs \emph{before authentication}, ahead of both the algorithm check
  and the signature check, so any well-formed token, forged or not, causes it,
  and pinning permitted algorithms does not prevent it.
\item It entered the library with the release that fixed an RSA-to-HMAC
  key-confusion vulnerability (CVE-2022-23541), and the key-type checks that
  implement that fix depend on it. The obvious performance optimisation
  therefore fails the library's own key-confusion tests.
\item A narrow fix preserves the defence and makes validation 4.15 times
  faster, but an attacker-chosen algorithm header can still reach the parse;
  only a pre-parsed key removes it for every input. In our sample, 94\% of
  call sites pin no algorithm and none pre-parses the key.
\end{itemize}
The magnitude of the cost is attributed causally to the bundled OpenSSL
release: the same OpenSSL calls issued from C, against eight identically built
releases, cost 394.4\,$\mu$s on 3.0.16 and 8.447\,$\mu$s on 3.5.8. It is
confirmed by cgroup CPU accounting in a containerised service on a single-node Kubernetes testbed. A comparison of seven
JWT libraries in four languages shows it follows a design choice rather than
JWT validation as such.

We believe the work suits \emph{Computers \& Security} because it concerns
the practical security engineering of authentication code. It shows how a
security hardening changes the cost profile of an authentication path, why
performance changes to verification code must be reviewed as security changes,
and what developers and library authors should do instead. It reports no new
attack. The defect has been reported to the library maintainers
(auth0/node-jsonwebtoken issue \#1046), and the public report describes the
constraint at the level of mechanism without an exploitation recipe. Every
number in the paper is generated from committed data in an open replication
package.

This manuscript is original, has not been published previously, and is not
under consideration for publication elsewhere. All authors have approved the
manuscript and agree with its submission. The authors declare no competing
interests. The use of generative AI in preparing the manuscript is declared in
the manuscript.

Thank you for considering our submission.

\medskip
\noindent
Yours sincerely,\\[4pt]
Tze Hui Liew, on behalf of all authors\\
Jannatul Ferdous Katha, Tasmia Tahmid Prova, Md Kishor Morol, Tze Hui Liew,
Dip Nandi

\end{document}
